\documentclass[12pt,a4paper]{article}

\usepackage[spanish]{babel}
\usepackage[utf8]{inputenc}
\usepackage[T1]{fontenc}
\usepackage{geometry}
\usepackage{listings}
\usepackage{xcolor}
\usepackage{titlesec}
\usepackage{hyperref}

\geometry{margin=2.5cm}

\title{
    \textbf{Documentación de implementación para la aplicación Pociones de Minecraft}
}

\author{}
\date{}

\definecolor{codegreen}{rgb}{0,0.6,0}
\definecolor{codegray}{rgb}{0.5,0.5,0.5}
\definecolor{codepurple}{rgb}{0.58,0,0.82}
\definecolor{backcolour}{rgb}{0.95,0.95,0.92}

\lstdefinestyle{mystyle}{
    backgroundcolor=\color{backcolour},
    commentstyle=\color{codegreen},
    keywordstyle=\color{blue},
    numberstyle=\tiny\color{codegray},
    stringstyle=\color{codepurple},
    basicstyle=\ttfamily\small,
    breaklines=true,
    captionpos=b,
    keepspaces=true,
    numbers=left,
    numbersep=5pt,
    showspaces=false,
    showstringspaces=false,
    showtabs=false,
    tabsize=2,
    language=Java
}

\lstset{style=mystyle}

\begin{document}

\maketitle

\tableofcontents

\newpage

\section{Introducción}

En este proyecto desarrollé la aplicación \textbf{Pociones de Minecraft}, utilizando una arquitectura cliente-servidor compuesta por Flutter para el frontend, Spring Boot para el backend y PostgreSQL como sistema gestor de base de datos.

El objetivo principal de la aplicación es permitir a los usuarios consultar y administrar recetas de pociones inspiradas en Minecraft, además de implementar funcionalidades como autenticación, comentarios y reacciones.

Durante el desarrollo implementé distintos patrones de diseño y buenas prácticas de programación orientada a objetos con el propósito de mantener una arquitectura organizada, escalable y reutilizable.

\section{Arquitectura general del proyecto}

Para organizar correctamente el sistema, dividí el proyecto en distintas capas:

\begin{itemize}
    \item Frontend Flutter
    \item Backend Spring Boot
    \item API REST
    \item Base de datos PostgreSQL
\end{itemize}

La arquitectura sigue el flujo:

\begin{center}
UI $\rightarrow$ Controladores $\rightarrow$ Servicios $\rightarrow$ Repositorios $\rightarrow$ Base de datos
\end{center}

Con esta separación logré desacoplar responsabilidades y facilitar el mantenimiento del sistema.

\section{Patrones de diseño utilizados}

\subsection{Patrón Singleton}

\subsubsection{¿Qué es?}

El patrón Singleton garantiza que únicamente exista una instancia de una clase durante toda la ejecución del programa.

\subsubsection{¿Cómo lo utilicé?}

Lo implementé dentro del servicio de autenticación en Flutter:

\begin{lstlisting}
class AuthService implements IAuthRepository {
  static final AuthService _instance = AuthService._internal();

  AuthService._internal();

  factory AuthService() {
    return _instance;
  }
}
\end{lstlisting}

\subsubsection{¿Para qué lo utilicé?}

Utilicé este patrón para:

\begin{itemize}
    \item Mantener una única sesión de autenticación
    \item Evitar múltiples instancias del cliente HTTP
    \item Compartir el token JWT en toda la aplicación
\end{itemize}

\subsection{Patrón Repository}

\subsubsection{¿Qué es?}

El patrón Repository permite encapsular el acceso a datos y separar la lógica de persistencia de la lógica de negocio.

\subsubsection{¿Cómo lo utilicé?}

Lo implementé utilizando Spring Data JPA:

\begin{lstlisting}
@Repository
public interface UserRepository extends JpaRepository<User, Long> {

    Optional<User> findByUsername(String username);

    Boolean existsByEmail(String email);
}
\end{lstlisting}

\subsubsection{¿Para qué lo utilicé?}

Con este patrón logré:

\begin{itemize}
    \item Abstraer consultas SQL
    \item Facilitar operaciones CRUD
    \item Mantener desacoplado el acceso a datos
\end{itemize}

\subsection{Patrón DTO (Data Transfer Object)}

\subsubsection{¿Qué es?}

El patrón DTO se utiliza para transferir información entre capas sin exponer directamente las entidades de la base de datos.

\subsubsection{¿Cómo lo utilicé?}

Lo implementé en las clases de autenticación:

\begin{lstlisting}
public class LoginRequest {

    private String username;

    private String password;
}
\end{lstlisting}

\begin{lstlisting}
public class JwtResponse {

    private String accessToken;

    private Long id;

    private String username;
}
\end{lstlisting}

\subsubsection{¿Para qué lo utilicé?}

Este patrón me permitió:

\begin{itemize}
    \item Validar la información recibida
    \item Proteger entidades internas
    \item Mejorar seguridad y organización
\end{itemize}

\subsection{Patrón Builder}

\subsubsection{¿Qué es?}

El patrón Builder facilita la creación de objetos complejos.

\subsubsection{¿Cómo lo utilicé?}

Lo implementé en la clase \texttt{UserDetailsImpl}:

\begin{lstlisting}
public static UserDetailsImpl build(User user) {

    List<GrantedAuthority> authorities =
        user.getRoles().stream()
            .map(role -> new SimpleGrantedAuthority(
                role.getName().name()))
            .collect(Collectors.toList());

    return new UserDetailsImpl(
            user.getId(),
            user.getUsername(),
            user.getEmail(),
            user.getPassword(),
            authorities);
}
\end{lstlisting}

\subsubsection{¿Para qué lo utilicé?}

Lo utilicé para:

\begin{itemize}
    \item Construir objetos de autenticación complejos
    \item Centralizar la creación de usuarios autenticados
    \item Mejorar reutilización del código
\end{itemize}

\subsection{Patrón Dependency Injection}

\subsubsection{¿Qué es?}

La inyección de dependencias permite que los objetos sean proporcionados automáticamente por el framework.

\subsubsection{¿Cómo lo utilicé?}

Implementé este patrón en controladores y servicios:

\begin{lstlisting}
@Autowired
AuthenticationManager authenticationManager;

@Autowired
UserRepository userRepository;
\end{lstlisting}

\subsubsection{¿Para qué lo utilicé?}

Gracias a este patrón pude:

\begin{itemize}
    \item Reducir acoplamiento
    \item Facilitar pruebas
    \item Mejorar mantenibilidad
\end{itemize}

\section{Buenas prácticas implementadas}

\subsection{Separación de responsabilidades}

Organicé el sistema para que cada clase tuviera una única responsabilidad:

\begin{itemize}
    \item Los Controllers manejan peticiones HTTP
    \item Los Services contienen lógica de negocio
    \item Los Repositories acceden a la base de datos
    \item Los Models representan entidades
\end{itemize}

\subsection{Uso de JWT}

Implementé autenticación mediante JWT.

\begin{lstlisting}
String jwt = jwtUtils.generateJwtToken(authentication);
\end{lstlisting}

Esto me permitió:

\begin{itemize}
    \item Utilizar autenticación stateless
    \item Mejorar seguridad
    \item Facilitar integración con Flutter
\end{itemize}

\subsection{Encriptación de contraseñas}

Protegí las contraseñas utilizando BCrypt.

\begin{lstlisting}
encoder.encode(signUpRequest.getPassword())
\end{lstlisting}

Con esto evité almacenar contraseñas en texto plano.

\subsection{Uso de HTTP Status Codes}

Utilicé códigos HTTP adecuados para representar estados del sistema:

\begin{itemize}
    \item 200 OK
    \item 201 Created
    \item 400 Bad Request
    \item 401 Unauthorized
\end{itemize}

\subsection{Uso de interfaces}

En Flutter implementé interfaces para desacoplar servicios:

\begin{lstlisting}
class AuthService implements IAuthRepository
\end{lstlisting}

Esto me permitió:

\begin{itemize}
    \item Mejorar escalabilidad
    \item Facilitar pruebas
    \item Permitir futuras implementaciones
\end{itemize}

\section{Seguridad implementada}

\subsection{Spring Security}

Utilicé Spring Security para proteger endpoints.

\begin{lstlisting}
.requestMatchers("/api/auth/**").permitAll()
.anyRequest().authenticated()
\end{lstlisting}

\subsection{Control de sesiones}

Configuré el sistema para trabajar con sesiones stateless.

\begin{lstlisting}
.sessionCreationPolicy(SessionCreationPolicy.STATELESS)
\end{lstlisting}

\subsection{Configuración CORS}

Configuré CORS para permitir comunicación segura entre Flutter Web y Spring Boot.

\begin{lstlisting}
configuration.setAllowedOrigins(List.of(
    "http://localhost:3000"
));
\end{lstlisting}

\section{Base de datos}

Utilicé PostgreSQL hospedado en Aiven como sistema gestor de base de datos.

Las principales entidades implementadas fueron:

\begin{itemize}
    \item User
    \item Role
    \item Potion
    \item PotionComment
    \item PotionReaction
\end{itemize}

También implementé relaciones:

\begin{itemize}
    \item OneToMany
    \item ManyToMany
    \item ManyToOne
\end{itemize}

\section{Despliegue}

Realicé el despliegue del backend utilizando Render.

Las tecnologías utilizadas fueron:

\begin{itemize}
    \item Spring Boot
    \item PostgreSQL
    \item Aiven
    \item Render
\end{itemize}

\section{Conclusión}

Durante el desarrollo de la aplicación Pociones de Minecraft implementé distintos patrones de diseño y buenas prácticas con el objetivo de construir una arquitectura organizada, reutilizable y escalable.

El uso conjunto de Flutter y Spring Boot me permitió desarrollar una aplicación moderna basada en APIs REST, autenticación JWT y separación adecuada de responsabilidades.
    
\end{document}